\documentclass[../../book-main.tex]{subfiles}
\graphicspath{{\subfix{../..}}}

\begin{document}

\chapter{Conditional Estimation and Generation}

In the previous chapters of this book, we have studied how to effectively and efficiently learn a representation for a variable $\x$ in the world with a distribution $p(\x)$ that has a low-dimensional support in a high-dimensional space. In general, a good representation or autoencoding should enable us to utilize the learned distribution of the data $\x$ and its representation $\z$ for various subsequent classification, 
estimation, and generation tasks under different conditions. The fundamental purpose of autoencoding is indeed to grant certain access to and \textit{control} over the distribution of the data. 

In the previous chapters, we have developed the methodology for representation learning and  autoencoding in a general, task-agnostic fashion. Here onwards, we will tailor our discussion
towards specific modern applications where we desire to learn useful
representations for specific structured data distributions, such as images, 3D
scenes, and videos, and where we have \textit{a priori} knowledge of certain
structured attributes of the data, such as partial observations, class labels, or text captions. In such cases, the modern practice is to establish such an autoencoding indirectly: rather than a direct encoder-decoder mapping between $\vx$ and $\vz$, we seek to establish a \textit{binding} or mapping between $\vx$ and $\vz$, which grants us the ability to generate samples of $\vx$ conditioned on
the (specified \textit{a priori}) attributes $\vz$, or even to generate these
attributes conditioned on a data sample $\vx$. This methodology is highly
related to the direct encoder-decoder autoencoding loops we have discussed in the previous Chapter,\footnote{And even uses it as a primitive in many cases, such as in the use of VAEs with conditional diffusion models.} but
it enables far more impressive applications, including all of the groundbreaking
multimodal capabilities of modern generative models.

In this chapter, we will discuss how to exploit such an autoencoding
representation for completion tasks in the data domain: a pedagogically-useful special case of
the general binding formalism. More precisely, we will study how to exploit the
(compressive) autoencoding for completing a data sample conditioned on a partial observation of the data. We will also study of data sampling and generation
conditioned on certain learned features or attributes of the data.
In this chapter and the next, we show how to apply these methods to learn the distribution and representation of many practical data which facilitate many real tasks in increasingly more general and challenging settings. 

In general, once we have learned the distribution of $\x$, it can be very useful for conducting subsequent inference in the world related to $\x$. As we will see, almost all practical applications can be converted to the problem that given an observation $\y$ that depends on $\x$, say
\begin{equation}
    \y = h(\x) + \boldsymbol{w},
\end{equation}
we need to solve the ``inverse problem'' of obtaining a good estimate $\hat \x(\y)$ for $\x$ or generating a sample $\hat{\x}$ that is consistent with the observation $\y \approx h(\hat{\x})$. To accomplish such tasks well, we essentially need to get a hold of the conditional distribution:
\begin{equation}
p(\x\mid \y) = \frac{p(\y\mid \x) p(\x) }{p(\y)}, 
\end{equation} 
where the last equation due to the Bayes rule.  Both the conditional estimation and conditional generation (or sampling or prediction) can be naturally inferred from the conditional distribution. 

In practice, however, initial information about the distributions of $\x$ and $\y$ and their relationship can be given in many different ways or forms. In general, they can mostly be categorized into four cases, which are, conceptually, increasingly more challenging:
\begin{itemize}
\item {\em Case 1:} Both a model for the distribution of $\x$ and the observation model $\y = h(\x) (+ \boldsymbol{w})$ is known, even in  an analytical form. This is typically the case for many classic signal processing problems, such as signal denoising, the sparse vector recovery problem we saw in Chapter \label{ch:classic} and the low-rank matrix recovery problem to be introduced below.
\item {\em Case 2:} We do not have a model for the distribution but only samples $\X = \{\x_1, \ldots, \x_N\}$ of $\x$, and the observation model $\y = h(\x) (+ \boldsymbol{w})$ is also known. A model for the distribution $p(\x)$ of $\x$ needs to be learned, and subsequently the conditional distribution $p(\x\mid \y)$. Natural image completion is one typical example for this class of problems.
\item {\em Case 3:} We only have the paired samples: $(\X, \Y) = \{ (\x_1, \y_1), \ldots, (\x_N, \y_N) \}$ of the two variables  $(\x, \y)$. The distributions of $\x$ and $\y$ and their relationship $h(\cdot)$ need to be learned from these paired sample data. For example, given many images and their captions, learning to conduct text-conditioned image generation is one such problem. 
\item {\em Case 4:} We only have the samples $\Y = \{\y_1, \ldots, \y_N\}$ of the observations $\y$, and the observation model $h(\cdot)$ needs to be known, at least in some parametric family $h(\cdot, \boldsymbol{\theta})$. The distribution $p(\x)$ and $p(\x \mid \y)$ need to be learned from $\hat \x$, estimated from $\Y$. For example, learning to render a new view from a sequence of calibrated or uncalibrated views is one such problem.
\end{itemize}
In this chapter, we will discuss general approaches to learn the desired distributions and solve the associated conditional estimation or generation for these cases, typically with a representative problem. 

\section{Conditional Estimation with a Known Data Distribution}
Notice that in the setting we have discussed in previous Chapters, the
autoencoding network is trained to reconstruct a set of complete samples of the random vector $\x$. This would allow us to generate samples of the learned (low-dimensional) distribution. In practice, however, the
low-dimensionality of the distribution is often exploited for
completion or prediction tasks. More specifically, we often need to
recover a sample $\x$ when parts of it are missing (or even corrupted). That is, we want to recover or predict the rest of $\x$
from observing only a fraction of it:
\begin{equation}
f: \mathcal{P}_{\Omega}(\x) \mapsto \hat{\x},
\end{equation}
where $\mathcal{P}_{\Omega}(\spcdot)$ represents a masking operation
(see \Cref{fig:matrix-completion} for an example).

In this section and the next, we will study the completion task under two different conditions: one is when the distribution of the data $\x$ is already given, either learned or known a priori (even in certain analytical form); one is when only raw samples of $\x$ are available and we need to learn the low-dimensional distribution from the samples in order to solve the completion task well.

\subsection{Low-rank Matrix Completion}
The low-rank {\em matrix
completion} problem is a classical problem for data completion when its distribution is low-dimensional and known. Consider a random vector $\x \in \mathbb{R}^m$
which is distributed according to a low-rank Gaussian (or
low-dimensional subspace) of rank $r < m$. Suppose we have drawn $n >
r$ independent samples from this distribution and form a matrix $\X =
[\x_1, \ldots, \x_n] \in \mathbb{R}^{m\times n}$. Obviously, the rank
of the matrix is
\begin{equation}
\mbox{rank}(\X) = r < \min\{m, n\}.
\end{equation}
Now let $\Omega$ indicate a set indices of observed entries of the
matrix $\X$. Let the observed entries to be:
\begin{equation}
\Y = \mathcal{P}_{\Omega}(\X).
\end{equation}
The remaining entries supported on $\Omega^c$ are unobserved or
missing. The problem is whether we can recover from $\Y$ the missing
entries correctly.  Figure \ref{fig:matrix-completion} shows one
example of matrix completion.
% \begin{figure}
%     \centering
%     \includegraphics[width=0.8\linewidth]{figures/Matrix-Completion.png}
%     \caption{An example for completing a matrix with correlations
% among entries.}
%     \label{fig:matrix-completion}
% \end{figure}

\begin{figure}
\centering
% \includegraphics[width=0.3\linewidth]{figures/Matrix-masked.png}
% \hspace{5mm}  \hspace{5mm}
% \includegraphics[width=0.3\linewidth]{figures/Matrix-completed.png}
\includegraphics[width=0.3\linewidth]{figs_chap5/masked-checkerboard.png}\;\;
\includegraphics[width=0.3\linewidth]{figs_chap5/mask.png}\;\;
\includegraphics[width=0.3\linewidth]{figs_chap5/checkerboard.png}
\caption{Illustration of  completing an image as low-rank matrix
  with some entries masked or corrupted. Left: the masked/corrupted
image $\Y$; middle: the mask $\Omega$; right: the completed image $\hat \X$.}
\label{fig:matrix-completion}
\end{figure}

Notice that the fundamental reason why such a matrix can be completed
is because columns and rows of the matrix are highly correlated and
they all lie on a low-dimensional subspace. For the example shown in
Figure \ref{fig:matrix-completion}, the dimension or the rank of the
matrix completed is only two. Hence the fundamental idea to recover
such a matrix is to seek a matrix that has the lowest rank among all
matrices that have entries agreeing with the observed ones:
\begin{equation}
\min \mbox{rank}(\X) \quad \mbox{subject to}
\quad
\Y = \mathcal{P}_{\Omega}(\X).
\label{eqn:rank-min}
\end{equation}
This is known as the {\em low-rank matrix completion} problem.
However, as the rank function is discontinuous and rank minimization
is in general an NP-hard problem, we often like to relax it with
something easier to optimize.

Based on our knowledge about compression from Chapter
\ref{ch:compression}, we could  promote the low-rankness of the
recovered matrix $\X$ by enforcing the lossy coding rate (or the
volume spanned by $\X$) of the data in $\X$ to be small:
\begin{equation}
\min R_\epsilon(\X) = \frac{1}{2} \log \det \left(\boldsymbol{I} +
\alpha  \X\X^\top \right) \quad \mbox{subject to}
\quad
\Y = \mathcal{P}_{\Omega}(\X).
\label{eqn:rate-min}
\end{equation}
The problem can viewed as a continuous relaxation of the above
low-rank matrix completion problem \eqref{eqn:rank-min} and it can be
solved via gradient descent. As one can show, the gradient descent
operator for the $\log\det$ objective is precisely minimizing a close
surrogate of the rank of the matrix $\X\X^\top$.

The rate distortion function is a nonlinear function, and its gradient
descent does not always guarantee finding the correct solution.
However, since the underlying structure sought for $\X$ is linear,
the rank function admits a rather effective convex relaxation: the
nuclear norm---the sum of all singular values of the matrix $\X$. As
shown in the compressive sensing literature, under fairly
broad conditions,\footnote{Typically, such conditions specify the
necessary and sufficient amount of entries needed for the completion
to be computationally feasible. These conditions have been
systematically characterized in \cite{Wright-Ma-2022}.} the matrix
completion problem \eqref{eqn:rank-min}
can be effectively solved by the following convex program:
\begin{equation}
\min \|\X\|_* \quad \mbox{subject to}
\quad
\Y = \mathcal{P}_{\Omega}(\X),
\label{eqn:nuclear-min}
\end{equation}
where the nuclear norm $\|\X\|_*$ is the sum of singular values of
$\X$. In practice, we often convert the above constrained convex optimization
program to an unconstrained one:
\begin{equation}
\min \|\X\|_*  + \lambda \|
\Y - \mathcal{P}_{\Omega}(\X)\|_F^2,
\label{eqn:nuclear-min-lagrangian}
\end{equation}
for some properly chosen $\lambda > 0$. Interested readers may refer to
\cite{Wright-Ma-2022} for how to develop algorithms  that can solve
the above programs  efficiently and effectively.
Figure \ref{fig:matrix-completion} shows a real example in which the
matrix $\hat \X$ is actually recovered by solving the above program.

It has been shown that for images (or more accurately textures) and 3D scenes
with low-rank structures, they can be very effectively completed via
solving optimization programs of the above kind, even if there is
additional corruption and distortion
\cite{Zhang2010TILTTI,Liang-ECCV2012,Yi_2023_ICCV}.

\section{Conditional Generation by Learning a Data Representation}

\subsection{Image Completion with Masked Auto-Encoding}
For a general image $\vX$ such as the one shown on the left of Figure
\ref{fig:crate_mae_pipeline}, we can no longer view it as a low-rank
matrix. However, humans still demonstrate remarkable ability to
complete a scene and recognize familiar objects despite severe
occlusion. This suggests that our brain has learned the
low-dimensional distribution of natural images and can use it for
completion, hence recognition. However, the distribution of all
natural images is not as simple as a low-dimensional linear subspace.
Hence a natural question is whether we can learn the more
sophisticated distribution of natural images and use it to perform
image completion?

\begin{figure}[t!]
%\vspace{-0.5in}
\begin{center}
  \includegraphics[width=0.99\textwidth]{figs_chap5/crate_mae_pipeline.png}
\end{center}
%\vspace{-0.1in}
\caption{\small \textbf{Diagram of the overall  (masked)
  autoencoding process.} The (image) token representations are
  transformed iteratively towards a parsimonious (e.g., compressed
  and sparse) representation by each encoder layer \(f^{\ell}\).
  Furthermore, such representations are transformed back to the
  original image by the decoder layers \(g^{\ell}\). Each encoder
  layer \(f^{\ell}\) is meant to be (partially) inverted by a
corresponding decoder layer \(g^{L - \ell}\).}
\label{fig:crate_mae_pipeline}
%\vspace{-0.2in}
\end{figure}

One natural approach is to find an encoding and decoding scheme by
solving the following {\em masked autoencoding} (MAE) program that
minimizes the reconstruction loss:
\begin{equation}\label{eq:mae_loss}
\min L_{\mathrm{MAE}}(f, g) \doteq \mathbb{E}\big[\norm{(g \circ
f)(\mathcal{P}_\Omega(\vX)) - \vX}_{2}^{2}].
\end{equation}
Unlike the matrix completion problem which has a simple underlying
structure, we should no longer expect that the encoding and decoding
mappings admit simple closed forms or the program can be solved by
explicit algorithms.

For a general natural image, we can no longer assume its columns or
rows are sampled from a low-dimensional subspace or a low-rank
Gaussian. However, it is reasonable to assume that the image consists
of multiple regions. Image patches in each region are similar and can
be modeled as one (low-rank) Gaussian or subspace. Hence, to exploit
the low-dimensionality of the distribution, the objective \textit{of the
encoder} $f$ is to transform $\X$ to a representation $\Z$:
\begin{equation}
    f: \X \mapsto \Z
\end{equation}
such that the distribution of $\Z$ can be well modeled as a mixture of subspaces, say $\{\vU_{[K]}\}$,
such that the overall coding rate and sparsity of the resulting
feature $\Z$ is the minimum:
\begin{equation}\label{eq:sparse_rr}
\min \mathbb{E}_{\vZ}[\Delta R(\vZ \mid \vU_{[K]}) - \lambda
\norm{\vZ}_{0}] = \mathbb{E}_{\vZ}[R(\vZ) - R^{c}(\vZ \mid
\vU_{[K]}) - \lambda \norm{\vZ}_{0}].
\end{equation}

As we have shown in the previous Chapter \ref{ch:representation}, the
encoder $f$ that minimizes the above objective can be constructed by
a sequence of transformer-like operators. As shown in the work of
\cite{Pai2024masked}, the decoder $g$ can be viewed and hence
constructed explicitly as the inverse process of the encoder $f$.
Figure \ref{fig:crate_mae_layers} illustrates the overall
architectures of both the encoder and the corresponding decoder at
each layer. The parameters of the encoder $f$ and decoder $g$ can be
learned by optimizing the reconstruction loss \eqref{eq:mae_loss} via
gradient descent.

\begin{figure}[t!]
%\vspace{-0.5in}
\centering
\includegraphics[width=0.99\textwidth]{figs_chap5/crate_mae_layers.png}
%\vspace{-0.1in}
\caption{\small \textbf{Diagram of each encoder layer
  (\textit{top}) and decoder layer (\textit{bottom}).} Notice that
  the two layers are highly anti-parallel --- each is constructed to
  do the operations of the other in reverse order. That is, in the
  decoder layer \(g^{\ell}\), the \(\ISTA\) block of \(f^{L - \ell}\)
  is partially inverted first using a linear layer, then the
  \(\MSSA\) block of \(f^{L - \ell}\) is reversed; this order
unravels the transformation done in \(f^{L - \ell}\).}
\label{fig:crate_mae_layers}
%\vspace{-0.2in}
\end{figure}

Figure \ref{fig:mae_autoencoding-small} shows some representative
results of the so-designed masked auto-encoder. More implementation details and
results of the masked autoencoder for natural image completion can be
found in Chapter \ref{ch:applications}.
\begin{figure}[t]
%\vspace{-0.5in}
\centering
\includegraphics[width=0.99\textwidth]{figs_chap5/crate_mae_autoencoding_small.pdf}
%\vspace{-0.1in}
\caption{\small \textbf{Autoencoding visualizations of CRATE-Base
  and ViT-MAE-Base \cite{he2022masked} with 75\% patches masked.}
  We observe that the reconstructions from CRATE-Base are on par
  with the reconstructions from ViT-MAE-Base, despite using \(<
  1/3\) of the parameters.
}
\label{fig:mae_autoencoding-small}
%\vspace{-0.2in}
\end{figure}

%\yima{Generalize low-rank matrix completion to masked autoencoding
% (MAE):  \href{https://arxiv.org/abs/2404.02446}{masked image completion}. }

%\yima{Maybe trying to optimize rate reduction type objective...
% Theoretical justification: maybe mention the recent work from Qing
% Qu's group about memorization and generalization (diffusion process
% for the PCA or mixture of low-dim Gaussians case). Probably some
% work on incremental optimal transport...}

%\section{Conditional Denoising and Diffusion}
%\label{sec:auto-denoising-diffusion}
%\yima{We now relax the requirement of sample-wise reconstruction to
% distribution-wise.}

The above (masked) autoencoding problems can be viewed to sample the learned low-dimensional distribution such that the generated samples are consistent with certain observations or conditions.

But let us examine the image completion task studied in the previous
section: Given the
visual part of an image $\X_v = \mathcal{P}_{\Omega}(\X)$, we try to
estimate the masked part $\X_m = \mathcal{P}_{\Omega^c}(\X)$. For realizations
$(\vXi_v, \vXi_m)$ of the random variable $\vX=(\vX_v, \vX_m)$, let
\[p_{\X_m \mid \X_v}(\vXi_m\mid \vXi_v)\]
be the conditional distribution of $\X_m$ given
$\X_v$. It is easy to show that the optimal solution to the  MAE
formulation \eqref{eq:mae_loss} is given by the conditional expectation:
\begin{equation}
  \argmin_{h = g \circ f}\, L_{\mathrm{MAE}}(h)
  % \doteq \mathbb{E}\big[\norm{(g \circ f)(\mathcal{P}_\Omega(\vX)) - \vX}_{2}^{2}].
  = \vXi_v \mapsto \vXi_v + \mathbb{E}[\X_m \mid \X_v=\vXi_v].
\end{equation}
In general, however, this expectation may not even be on the
low-dimensional distribution of natural images! This partially
explains why some of the recovered patches in \Cref{fig:mae_autoencoding-small}
are a little blurry.

For many practical purposes, we would like to learn (a representation
of) the conditional distribution $p_{\X_m \mid \X_v}$, or equivalently
$p_{\X \mid \X_v}$,
and then sample this distribution directly. Notice that, when the
distribution of $\X$ is low-dimensional, it is possible that if a
sufficient part of $\X$, $\X_v$, is observed, it fully determines
$\X$ hence the missing part $\X_m$. In other words, the distribution
$p_{\X \mid \X_v}$ is a generalized function---if $\vX$ is fully determined by $\vX_v$ it is the delta function, and more generally one of its exotic cousins.

\subsection{Conditional Sampling with Measurement Matching}
\label{sec:conditioned-decoding}
To generalize the above (image) completion problems, we may consider
that a random vector $\vx \sim p$ is partially observed through a more
general observation function:
\begin{equation}\label{eq:measurement-matching-observation}
\vy = h(\vx) + \vw,
\end{equation}
where $\vw$ usually stands for some random measurement noise, say of
a Gaussian distribution $\vw \sim \mathcal{N}(\mathbf{0}, \sigma^2
\boldsymbol{I})$. It is easy to see that, for so related $\x$ and $\y$, their joint distribution $p(\x, \y)$ is naturally nearly degenerate if the noise $\vw$ is small. To a large extent, we may view $p(\x, \y)$ as a noisy version of a hypersurface defined by the function $\y = h(\x)$ in the joint space $(\x, \y)$.  Practically speaking, we will consider a setting more akin to masked autoencoding than to pure matrix completion, where we always have access to a corresponding clean sample $\vx$ for every observation $\vy$ we receive.\footnote{In some more specialized applications, in particular in scientific imaging, it is of interest to be able to learn to generate samples
from the posterior $p_{\x \mid \y}$ without access to any clean/ground-truth samples of $\x$. We give a brief overview of methods for this setting in the
end-of-chapter notes.}

Like image/matrix completion, we
are often faced with a setting where $\vy$ denotes a degraded or otherwise
``lossy''
observation of the input $\vx$. This can manifest in quite different
forms. For example, in various scientific or
medical imaging problems, the measured data $\vy$ may be a compressed and
corrupted observation of the underlying data $\vx$; whereas in 3D vision tasks,
$\vy$ may represent an image captured by a camera of a physical object with an
unknown (low-dimensional) pose $\vx$.
Generally, by virtue of mathematical modeling (and, in some cases, co-design
of the measurement system), we know $h$ and can evaluate it on any input, and
we can exploit this knowledge to help reconstruct and sample $\vx$.

\begin{figure}[t]
  \centering
  \begin{subfigure}{0.45\textwidth}
    \vspace{0.75cm}
    \centering
    \begin{tikzcd}[column sep=1.5cm]
      \vx \arrow[r, "h(\vx) + \vw"] & \vy \arrow[r, "p_{\vx \mid \vy}"]
      & \posteriorsample \arrow[r, "\posteriorsample + t\vg"] & \posteriorsample_t
    \end{tikzcd}
    \vspace{0.75cm}
    \caption{}
    \label{fig:left}
  \end{subfigure}
  \hfill
  \begin{subfigure}{0.45\textwidth}
    \centering
    \begin{tikzcd}[column sep=1.5cm, row sep=0.5cm]
    & \vy \\
    \vx \arrow[ur, "h(\vx) + \vw"] \arrow[dr, "\vx + t\vg"] & \\
    & \vx_t
    \end{tikzcd}
    \caption{}
    \label{fig:right}
  \end{subfigure}
  \caption{Statistical dependency diagrams for the conditional sampling process.
  \textbf{Left:} In a direct (conceptual) application of the diffusion-denoising scheme
  we have developed in \Cref{ch:compression} to conditional sampling, we use
  samples from the posterior $p_{\vx \mid \vy}$ to train denoisers directly on
  the posterior at different noise levels, then use them to generate new
  samples. In practice, however, we do not normally have direct samples from the
  posterior, but rather paired samples $(\vx, \vy)$ from the joint.
  \textbf{Right:} It turns out that it suffices to have only noisy observations
  of $\vx$ to realize the denoisers corresponding to $p_{\posteriorsample_t \mid
  \posteriorsample}$: this follows from conditional independence of $\vx_t$ and
  $\vy$ given $\vx$. It implies that $p_{\posteriorsample_t} = p_{\vx_t \mid
  \vy}$, which gives a score function for denoising that consists of the
  unconditional score function, plus a correction term that enforces measurement
  consistency.}
  \label{fig:posterior-sampling-cds}
\end{figure}


At a technical level, we want the learned representation of the data
to facilitate us to sample the conditional distribution $p_{\vx\mid
\vy}$, also known as the posterior, effectively and efficiently. More precisely,
write $\vnu$ to denote a realization of the random variable $\vy$.
We want to generate samples $\hat{\x}$ such that:
\begin{equation}\label{eq:goal-sample-posterior}
  \hat{\x} \sim   p_{\vx \mid \y}(\spcdot \mid \vy=\vnu).
\end{equation}

Recall that in \Cref{sub:compression_denoising}, we have developed a natural and
effective way to produce \textit{unconditional} samples of the data distribution
$p$. The ingredients are the denoisers $\bar{\x}^\ast(t, \vxi) = \bE[ \x \mid
\vx_t=\vxi ]$, or their learned approximations $\bar{\x}_{\theta}(t, \vxi)$,
for different levels of noisy observations $\x_t = \x + t \vg$ (and $\vxi$ for
their realizations) under
Gaussian noise $\vg \sim \cN(\mathbf{0}, \vI)$, and $t \in
[0, T]$ with a choice of times $0 = t_1 < \hdots < t_{L} = T$ at which to
perform the iterative denoising, starting from $\hat{\vx}_{t_L} \sim
\cN(\mathbf{0}, T^2 \vI)$ (recall \Cref{eq:denoising-iteration-basic}).\footnote{Recall from our discussion in \Cref{sub:sampling_denoising}
that a few small improvements to this basic iterative denoising scheme are
sufficient to bring competitive practical performance. For clarity as we develop
conditional sampling, we will focus here on the simplest instantiation.}
We could directly apply this scheme to generate samples from the posterior
$p_{\vx \mid \vy}$ \textit{if} we had access to a dataset of samples
$\posteriorsample \sim p_{\vx \mid \y}(\spcdot \mid \vnu)$ for each realization
$\vnu$ of $\vy$, by generating noisy observations
$\posteriorsample_t$ and training denoisers to approximate $\bE[
  \posteriorsample \mid \posteriorsample_t=\spcdot, \vy=\vnu ]$, the mean of the posterior under
the noisy observation (see \Cref{fig:posterior-sampling-cds}(a)).
However, performing this resampling given only paired samples $(\vx, \vy)$ from the
joint distribution (say by binning the samples over values of $\vy$) requires
prohibitively many samples for high-dimensional data, and alternate approaches
explicitly or implicitly rely on density estimation, which similarly suffers from the
curse of dimensionality.

Fortunately, it turns out that this is not necessary.
Consider the alternate statistical dependency diagram in
\Cref{fig:posterior-sampling-cds}(b), which corresponds to the random variables
in the usual denoising-diffusion process, together with the measurement $\vy$. 
Because our assumed observation model \eqref{eq:measurement-matching-observation}
implies that $\vx_t$ and $\vy$ are independent conditioned on $\vx$, we have for
any realization $\vnu$ of $\vy$
\begin{align*}
  p_{\posteriorsample_t\mid \vy}(\spcdot \mid \vnu)
  &= \int
  \underbrace{p_{\posteriorsample_t \mid \posteriorsample}(\spcdot\mid \vxi)}_{=
  \cN(\vxi, t^2 \vI)}
  \spcdot \underbrace{p_{\posteriorsample \mid \vy}}_{= p_{\vx\mid \vy}}(\vxi \mid
  \vnu)
  \odif \vxi
  \\
  &=
  \int p_{\vx_t \mid \vx, \vy}(\spcdot\mid \vxi, \vnu) \spcdot p_{\vx \mid
  \vy}(\vxi\mid \vnu) \odif \vxi
  \\
  &=
  \int p_{\vx_t, \vx \mid \vy}(\spcdot, \vxi \mid \vnu) \odif \vxi
  \\
  &= p_{\vx_t \mid \vy}(\spcdot\mid \vnu).
\end{align*}
Above, the first line recognizes an equivalence between the distributions
arising in \Cref{fig:posterior-sampling-cds}(a, b); the second line applies this
together with conditional independence of $\vx_t$ and $\vy$ given $\vx$; the
third line uses the definition of conditional probability; and the final line
marginalizes over $\vx$.
Thus, the denoisers from the conceptual posterior sampling process are equal to
$\bE[\vx \mid \vx_t=\spcdot, \vy=\vnu]$, which we can learn solely from paired samples $(\vx,
\vy)$,
% and the posterior with respect to $\vx_t$ and $\vy$ can be expressed using
% Bayes' rule as
% \begin{equation}
%   p(\vx \mid \vx_t, \vy) = \frac{p(\vx_t \mid \vx) p(\vx \mid \vy)}{p(\vx_t \mid
%   \vy)},
% \end{equation}
% where we used that $\vx_t$ is an independent Gaussian conditional on $\vx$.
and by Tweedie's formula (\Cref{thm:tweedie}), we can express these denoisers in
terms of the score function of $p_{\vx_t \mid \vy}$, which, by Bayes' rule,
satisfies
\begin{equation}
  p_{\vx_t \mid \vy}(\vxi \mid \vnu) 
  = \frac{p_{\vy \mid \vx_t}(\vnu \mid \vxi) p_{\vx_t}(\vxi)}{p_{\vy}(\vnu)}.
\end{equation}
Recall that the density of $\vx_t$ is 
given by $p_t = \varphi_{t} \ast p$, where $\varphi_{t}$ denotes the standard
Gaussian density with zero mean and covariance $t^2 \vI$ and $\ast$ denotes
convolution. This is nothing but the \textit{unconditional score function}
obtained from the standard diffusion training that we developed in
\Cref{sub:compression_denoising}!
The conditional score function then satisfies, for any realization $(\vxi,
\vnu)$ of $(\vx_t, \vy)$,
\begin{equation}
  \nabla_{\vxi} \log p_{\vx_t \mid \vy}(\vxi\mid \vnu)
  =
  \underbrace{\nabla_{\vxi} \log p_t(\vxi)}_{\text{score matching}}
  +
  \underbrace{\nabla_{\vxi} \log p_{\vy \mid \vx_t}(\vnu \mid \vxi)}_{\text{measurement matching}},
\end{equation}
giving (by Tweedie's formula) our proposed denoisers as
\begin{align*}
  \bE[ \vx \mid \vx_t=\vxi, \vy=\vnu]
  &=
  \vxi + t^2 \nabla_{\vxi}\log p_t(\vxi) + t^2 \nabla_{\vxi}\log p_{\vy \mid
  \vx_t}(\vnu \mid \vxi)
  \\
  &=
  \bE[\vx \mid \vx_t=\vxi] + t^2 \nabla_{\vxi}\log p_{\vy \mid \vx_t}(\vnu\mid
  \vxi).
  \labelthis \label{eq:posterior-sampling-denoiser-decomposition}
\end{align*}
The resulting operators are interpretable as a \textit{corrected} version of the
unconditional denoiser for the noisy observation, where the correction term (the
so-called ``measurement matching'' term) enforces consistency with the
observations $\vy$. The reader should take care to note to which argument the
gradient operators are applying in the above score functions in order to fully
grasp the meaning of this operator.

The key remaining issue in making this procedure computational is to prescribe
how to compute the measurement matching correction, since in general we do not
have a closed-form
expression for the likelihood $p_{\vy \mid \vx_t}$ except for when $t
= 0$. Before taking up this problem, we discuss an illustrative concrete example
of the entire process, continuing from those we have developed in
\Cref{sub:compression_denoising}.

\begin{example}\label{example:denoising-conditional-gaussian}
  Consider the case where the data distribution is Gaussian with mean $\vmu \in
  \bR^D$ and
  covariance $\vSigma \in \bR^{D \times D}$, i.e., $\vx \sim \cN(\vmu, \vSigma)$.
  Assume that $\vSigma \succeq \Zero$ is nonzero.
  Moreover, in the measurement model \eqref{eq:measurement-matching-observation},
  suppose we obtain linear measurements of $\vx$ with independent Gaussian
  noise, where $\vA \in \bR^{d \times D}$ and $\vy = \vA \vx + \sigma \vw$
  with $\vw \sim \cN(\Zero, \vI)$ independent of $\vx$.
  Then $\vx =_{d} \vSigma^{1/2} \vg + \vmu$, where $\vg \sim \cN(\Zero, \vI)$ is
  independent of $\vw$ and $\vSigma^{1/2}$ is the unique positive square root of
  the covariance matrix $\vSigma$, and
  after some algebra, we can then write
  \begin{equation*}
    \begin{bmatrix}
      \vx \\
      \vy
    \end{bmatrix}
    =_{d}
    \begin{bmatrix}
      \vSigma^{1/2} & \Zero \\
      \vA \vSigma^{1/2} & \sigma \vI
    \end{bmatrix}
    \begin{bmatrix}
      \vg \\
      \vw
    \end{bmatrix}
    +
    \begin{bmatrix}
      \vmu \\
      \vA\vmu
    \end{bmatrix}.
  \end{equation*}
  By independence, we have that $(\vg, \vw)$ is jointly Gaussian, which means that
  $(\vx, \vy)$ is also jointly Gaussian, as the affine image of a jointly Gaussian
  vector.
  Its covariance matrix is given by
  \begin{equation*}
    \begin{bmatrix}
      \vSigma^{1/2} & \Zero \\
      \vA \vSigma^{1/2} & \sigma \vI
    \end{bmatrix}
    \begin{bmatrix}
      \vSigma^{1/2} & \Zero \\
      \vA \vSigma^{1/2} & \sigma \vI
    \end{bmatrix}^\top
    =
    \begin{bmatrix}
      \vSigma & \vSigma \vA^\top \\
      \vA \vSigma & \vA \vSigma \vA^\top + \sigma^2 \vI
    \end{bmatrix}.
  \end{equation*}
  Now, we apply the fact that conditioning a random vector with joint Gaussian
  distribution on a subset of coordinates is again a Gaussian distribution
  (\Cref{exercise:conditional_gaussian}).
  By this, we obtain that
  \begin{equation}
    p_{\x \mid \y}(\spcdot \mid \vnu) = \cN\left(
    \underbrace{%
      \vmu + \vSigma\vA^\top \left(\vA\vSigma\vA^\top + \sigma^2 \vI\right)^{-1} 
      (\vnu - \vA\vmu)}_{\vmu_{\vx \mid \vy}(\vnu)},
      \underbrace{%
      \vSigma - \vSigma \vA^\top \left(\vA\vSigma\vA^\top + \sigma^2
      \vI\right)^{-1} \vA \vSigma}_{\vSigma_{\vx\mid \vy}}
    \right).
  \end{equation}
  By the equivalence we have derived above, we get by another application of
  \Cref{exercise:conditional_gaussian}
  \begin{equation}\label{eq:conditional-posterior-denoiser-gaussian-case}
    \bE[ \vx \mid \vx_t=\vxi, \vy=\vnu ]
    =
    \vmu_{\vx\mid \vy}(\vnu) + \vSigma_{\vx\mid\vy}\left(
    \vSigma_{\vx\mid \vy} + t^2 \vI 
    \right)^{-1}\left(
    \vxi - \vmu_{\vx\mid \vy}(\vnu)
    \right).
  \end{equation}
  The functional form of this denoiser is quite simple, but it carries an
  unwieldy dependence on the problem data $\vmu$, $\vSigma$, $\vA$, and
  $\sigma^2$.
  We can gain further insight into its behavior by comparing it with
  \Cref{eq:posterior-sampling-denoiser-decomposition}.
  We have as usual
  \begin{equation}\label{eq:posterior-denoiser-gaussian-case}
    \bE[\vx \mid \vx_t=\vxi]
    = \vmu + \vSigma\left(\vSigma + t^2 \vI\right)^{-1} (\vxi - \vmu),
  \end{equation}
  which is rather simple---suggesting that the measurement matching term is
  rather complicated. To confirm this, we can calculate the likelihood $p_{\vy
  \mid \vx_t}$ directly using the following expression for the joint
  distribution of $(\vx_t, \vy)$:
  \begin{equation}
    \begin{bmatrix}
      \vx \\
      \vx_t \\
      \vy
    \end{bmatrix}
    =_{d}
    \begin{bmatrix}
      \vSigma^{1/2} & \Zero & \Zero \\
      \vSigma^{1/2} & t\vI & \Zero \\
      \vA \vSigma^{1/2} & \Zero & \sigma \vI
    \end{bmatrix}
    \begin{bmatrix}
      \vg \\
      \vg' \\
      \vw
    \end{bmatrix}
    +
    \begin{bmatrix}
      \vmu \\
      \vmu \\
      \vA\vmu
    \end{bmatrix},
  \end{equation}
  where $\vg' \sim \cN(\Zero, \vI)$ independent of the other Gaussians.
  This is again a jointly Gaussian distribution; restricting to only the final
  two rows, we have the covariance
  \begin{equation*}
    \begin{bmatrix}
      \vSigma^{1/2} & t\vI & \Zero \\
      \vA \vSigma^{1/2} & \Zero & \sigma \vI
    \end{bmatrix}
    \begin{bmatrix}
      \vSigma^{1/2} & t\vI & \Zero \\
      \vA \vSigma^{1/2} & \Zero & \sigma \vI
    \end{bmatrix}^\top
    =
    \begin{bmatrix}
      \vSigma + t^2\vI & \vSigma\vA^\top \\
      \vA\vSigma & \vA\vSigma\vA^\top + \sigma^2 \vI
    \end{bmatrix}.
  \end{equation*}
  Another application of \Cref{exercise:conditional_gaussian} then gives us
  \begin{equation}
    p_{\y \mid \vx_t}(\spcdot \mid \vxi) = \cN\left(
    \underbrace{%
      \vA\vmu + \vA\vSigma\left( \vSigma + t^2 \vI \right)^{-1}
      \left(
        \vxi - \vmu
      \right)
      }_{\vmu_{\vy \mid \vx_t}(\vxi)},
      \underbrace{%
        \vA\vSigma\vA^\top + \sigma^2 \vI - \vA\vSigma \left( \vSigma
        + t^2 \vI\right)^{-1} \vSigma \vA^\top
      }_{\vSigma_{\vy\mid \vx_t}}
    \right).
  \end{equation}
  Now notice that $\vmu_{\vy \mid \vx_t}(\vxi) = \vA \bE[\vx \mid \vx_t=\vxi]$.
  So, by the chain rule,
  \begin{align*}
    t^2 \nabla_{\vxi} \log p_{\vy \mid \vx_t}(\vnu \mid \vxi)
    &=
    t^2 \nabla_{\vxi}\left[
      -\frac{1}{2}
      (\vnu - \vA \bE[\vx \mid \vx_t=\vxi])^\top
      \vSigma_{\vy \mid \vx_t}^{-1}
      (\vnu - \vA \bE[\vx \mid \vx_t=\vxi])
      \right]
    \\
    &= t^2 (\vSigma+t^2 \vI)^{-1}\vSigma\vA^\top\vSigma_{\vy \mid \vx_t}^{-1}\left(
    \vnu - \vA \bE[\vx \mid \vx_t=\vxi] \right).
    \labelthis
    \label{eq:conditional-posterior-measurementmatching-gaussian-case}
  \end{align*}
  This gives us a more interpretable decomposition of the conditional posterior
  denoiser \eqref{eq:conditional-posterior-denoiser-gaussian-case}: following
  \Cref{eq:posterior-sampling-denoiser-decomposition}, it is the sum of the
  unconditional posterior denoiser \eqref{eq:posterior-denoiser-gaussian-case}
  and the measurement matching term
  \eqref{eq:conditional-posterior-measurementmatching-gaussian-case}.
  We can further analyze the measurement matching term. Notice that
  \begin{equation}
    \vSigma_{\vy \mid \vx_t}
    =
    \sigma^2 \vI + \vA\vSigma^{1/2} \left(
      \vI - \vSigma^{1/2}\left(
      \vSigma + t^2 \vI
      \right)^{-1}
      \vSigma^{1/2}
    \right)\vSigma\vA^\top.
  \end{equation}
  If we let $\vSigma = \vV \vLambda \vV^\top$ denote an eigenvalue decomposition
  of $\vSigma$, where $(\vv_i)$ are the columns of $\vV$, we can further write
  \begin{align}
    \vSigma^{1/2}\left(
    \vI - \vSigma^{1/2}\left(
    \vSigma + t^2 \vI
    \right)^{-1}
    \vSigma^{1/2}
    \right) \vSigma^{1/2}
    &=
    t^2 \vV \vLambda^{1/2}\left(
      \vLambda + t^2 \vI
    \right)^{-1}
    \vLambda^{1/2}
    \vV^\top
    \\
    &=
    t^2 \sum_{i=1}^D
    \frac{\lambda_i}{\lambda_i + t^2}
    \vv_i \vv_i\adj.
  \end{align}
  Then for any eigenvalue of $\vSigma$ equal to zero, the corresponding summand
  is zero; and
  writing $\lambda_{\min}(\vSigma)$ for the smallest positive eigenvalue of
  $\vSigma$ (it has at least one positive eigenvalue, by assumption), we have
  (in a sense that can be made quantitatively precise) that whenever $t \ll
  \sqrt{\lambda_{\min}(\vSigma)}$, it holds
  \begin{equation}
    \frac{\lambda_i t^2}{\lambda_i + t^2} \approx 0.
  \end{equation}
  So, when $t \ll \sqrt{\lambda_{\min}(\vSigma)}$, we have the approximation
  \begin{equation}
    \vSigma_{\vy \mid \vx_t} \approx \sigma^2 \vI.
  \end{equation}
  The righthand side of this approximation is equal to $\vSigma_{\vy \mid \vx}$.
  So we have in turn
  \begin{equation}
    \nabla_{\vxi} \log p_{\vy \mid \vx_t}(\vnu \mid \vxi)
    \approx
    \nabla_{\vxi} \log p_{\vy \mid \vx}(\vnu \mid \bE[\vx \mid \vx_t = \vxi]).
    \label{eq:conditional-posterior-measurementmatching-gaussian-case-dps-approx}
  \end{equation}

  \begin{figure}[tbp]
    \centering
    % Top row (3 images)
    \begin{subfigure}{0.32\textwidth}
      \includegraphics[width=\linewidth]{figs_chap5/samples_step_000_t_100_largenoise.png}
      % \caption{}
      % \label{}
    \end{subfigure}
    \hfill
    \begin{subfigure}{0.32\textwidth}
      \includegraphics[width=\linewidth]{figs_chap5/samples_step_050_t_50_largenoise.png}
      % \caption{}
      % \label{}
    \end{subfigure}
    \hfill
    \begin{subfigure}{0.32\textwidth}
      \includegraphics[width=\linewidth]{figs_chap5/samples_step_100_t_0_largenoise.png}
      % \caption{}
      % \label{}
    \end{subfigure}

    \vspace{2mm} % Add some vertical space between rows

    % Bottom row (3 images)
    \begin{subfigure}{0.32\textwidth}
      \includegraphics[width=\linewidth]{figs_chap5/samples_step_000_t_100_smallnoise.png}
      % \caption{}
      % \label{}
    \end{subfigure}
    \hfill
    \begin{subfigure}{0.32\textwidth}
      \includegraphics[width=\linewidth]{figs_chap5/samples_step_010_t_90_smallnoise.png}
      % \caption{}
      % \label{}
    \end{subfigure}
    \hfill
    \begin{subfigure}{0.32\textwidth}
      \includegraphics[width=\linewidth]{figs_chap5/samples_step_100_t_0_smallnoise.png}
      % \caption{}
      % \label{}
    \end{subfigure}

    \caption{Numerical simulation of the conditional sampling setup
    \eqref{eq:measurement-matching-observation}, with Gaussian data, linear
    measurements, and Gaussian noise. We simulate $D=2$ and $d=1$, with $\vSigma
    = \ve_1 \ve_1^\top + \tfrac{1}{4} \ve_2 \ve_2^\top$, $\vmu = \Zero$, and
    $\vA = \ve_1^\top$. The underlying signal $\vx$ is marked with a black star,
    and the measurement $\vy$ is marked with a black circle.
    Each individual plot corresponds to a different value of sampler time
    $t_{\ell}$, with different rows corresponding to different observation noise
    levels $\sigma^2$.
    In each plot, the covariance matrix of $\vx$ is plotted in gray, the
    posterior covariance matrix and posterior mean of $p_{\vx \mid \vy}$ are
    plotted in blue (with the posterior mean marked by a blue ``x''), and
    contours for $p_{\vx_{t_{\ell}} \mid \vy}$ are drawn in green. The sampler
    hyperparameters are $T=1$, $L=100$, and we draw $100$
    independent samples to initialize the samplers. Samplers are implemented
    with the closed-form denoisers derived in
    \Cref{example:denoising-conditional-gaussian}, with those using the
    approximation
    \eqref{eq:conditional-posterior-measurementmatching-gaussian-case-dps-approx}
    marked with red triangles, and those using the exact conditional posterior
    denoiser marked with blue circles.
    \textbf{Top:} For large observation noise $\sigma = 0.5$, both the exact
    conditional posterior denoiser and the approximate one do a good job of
    converging to the posterior $p_{\vx \mid \vy}$. Sampling time (corresponding
    to time in the ``forward process'', so larger times mean larger noise)
    decreases from left to right.  The convergence dynamics for the exact and
    approximate measurement matching term are similar. \textbf{Bottom:} For
    smaller observation noise $\sigma = 0.1$, the
    approximate measurement matching term leads to extreme bias in the sampler
    (red triangles): samples rapidly converge to an affine subspace
    of points that are consistent, modulo some shrinkage from the posterior mean
    denoiser, with the measured ground truth, and later sampling iterations are
    unable to recover the lost posterior variance along this dimension. Note
    that different times $t_{\ell}$ are plotted in the bottom row, compared to
    the top row, to show the rapid collapse of the approximation to the
    posterior along the measurement dimension.}
    \label{fig:conditional_sampling_computational_gaussian}
  \end{figure}

  \Cref{eq:conditional-posterior-measurementmatching-gaussian-case-dps-approx}
  is, of course, a direct consequence of our calculations above. However, notice
  that if we directly interpret this approximation, it is \textit{ab initio}
  tractable: the likelihood $p_{\vy \mid \vx} = \cN(\vA\vx, \sigma^2 \vI)$ is
  a simple Gaussian distribution centered at the observation, and the
  approximation to the measurement matching term that we arrive at can be
  interpreted as simply evaluating the log-likelihood at the conditional
  expectation $\bE[\vx \mid \vx_t = \vxi]$, then taking gradients with respect
  to $\vxi$ (and backpropagating through the conditional expectation, which is
  given here by \Cref{eq:posterior-denoiser-gaussian-case}).
  Nevertheless, note that the approximation in
  \Cref{eq:conditional-posterior-measurementmatching-gaussian-case-dps-approx}
  requires $t \ll \sqrt{\lambda_{\min}(\vSigma)}$, and that it is never accurate
  in general when this condition does not hold, even in this Gaussian setting.

  To gain insight into the effect of the convenient approximation
  \eqref{eq:conditional-posterior-measurementmatching-gaussian-case-dps-approx},
  we implement and simulate a simple numerical experiment in the Gaussian
  setting in \Cref{fig:conditional_sampling_computational_gaussian}. 
  The sampler we implement is a direct implementation of the simple scheme
  \eqref{eq:denoising-iteration-basic} we have developed in \Cref{ch:compression}
  and recalled above, using the true conditional posterior denoiser, i.e.\
  \Cref{eq:conditional-posterior-denoiser-gaussian-case} (top row of
  \Cref{fig:conditional_sampling_computational_gaussian}), and the 
  convenient approximation to this denoiser made with the decomposition
  \eqref{eq:posterior-sampling-denoiser-decomposition}, the posterior denoiser
  \eqref{eq:posterior-denoiser-gaussian-case}, and the measurement
  matching approximation
  \eqref{eq:conditional-posterior-measurementmatching-gaussian-case-dps-approx}
  (bottom row of \Cref{fig:conditional_sampling_computational_gaussian}).
  We see that even in the simple Gaussian setting, the approximation to the
  measurement matching term we have made is not without its
  drawbacks---specifically, at small noise levels $\sigma^2 \ll 1$, it leads to
  rapid collapse of the variance of the sampling distribution along directions
  that are parallel to the rows of the linear measurement operator $\vA$, which
  cannot be corrected by later iterations of sampling. We can intuit this from
  the approximation
  \eqref{eq:conditional-posterior-measurementmatching-gaussian-case-dps-approx}
  and the definition of the denoising iteration
  \eqref{eq:denoising-iteration-basic}, given
  \Cref{eq:posterior-sampling-denoiser-decomposition}: for $\sigma^2 \ll 1$,
  early steps of sampling effectively take gradient descent steps with a very
  large step size on the likelihood, via
  \Cref{eq:conditional-posterior-measurementmatching-gaussian-case-dps-approx},
  which leads the sampling distribution to get ``stuck'' in a collapsed state.

  % On the other hand, the second (more tractable) scheme in
  % \Cref{fig:posterior-sampling-cds}, which involves generating $\vx_t$ by adding
  % noise to $\vx$ and then conditioning on $\vy$ after, produces a different
  % joint distribution
  % \begin{equation}
  %   \begin{bmatrix}
  %     \vx \\
  %     \vx_t \\
  %     \vy
  %   \end{bmatrix}
  %   =_{d}
  %   \begin{bmatrix}
  %     \vSigma^{1/2} & \Zero & \Zero \\
  %     \vSigma^{1/2} & t\vI & \Zero \\
  %     \vA \vSigma^{1/2} & \Zero & \sigma \vI
  %   \end{bmatrix}
  %   \begin{bmatrix}
  %     \vg \\
  %     \vg' \\
  %     \vw
  %   \end{bmatrix}
  %   +
  %   \begin{bmatrix}
  %     \vmu \\
  %     \vmu \\
  %     \vA\vmu
  %   \end{bmatrix},
  % \end{equation}
  % where $\vg' \sim \cN(\Zero, \vI)$ independent of the other Gaussians.
  % This is again a jointly Gaussian distribution, with covariance
  % \begin{equation*}
  %   \begin{bmatrix}
  %     \vSigma^{1/2} & \Zero & \Zero \\
  %     \vSigma^{1/2} & t\vI & \Zero \\
  %     \vA \vSigma^{1/2} & \Zero & \sigma \vI
  %   \end{bmatrix}
  %   \begin{bmatrix}
  %     \vSigma^{1/2} & \Zero & \Zero \\
  %     \vSigma^{1/2} & t\vI & \Zero \\
  %     \vA \vSigma^{1/2} & \Zero & \sigma \vI
  %   \end{bmatrix}^\top
  %   =
  %   \begin{bmatrix}
  %     \vSigma & \vSigma & \vSigma\vA^\top \\
  %     \vSigma & \vSigma + t^2\vI & \vSigma\vA^\top \\
  %     \vA \vSigma & \vA\vSigma & \vA\vSigma\vA^\top + \sigma^2 \vI
  %   \end{bmatrix}.
  % \end{equation*}
  % Using \Cref{exercise:conditional_gaussian} once again, we derive the posterior
  % as
  % \begin{equation}
  %   p_{\vx \mid \vy, \vx_t}
  %   =
  %   \cN\left(
  %     ...
  %   \right)
  % \end{equation}
\end{example}

\Cref{example:denoising-conditional-gaussian} suggests a convenient
approximation for the measurement matching term
\eqref{eq:conditional-posterior-measurementmatching-gaussian-case-dps-approx},
which can be made beyond the Gaussian setting of the example. To
motivate this approximation in greater generality, notice that
by conditional independence of $\vy$ and $\vx_t$ given $\vx$, we can write
\begin{equation}
  p_{\vy \mid \vx_t}(\vnu \mid \vxi)
  =
  \int p_{\vy \mid \vx}(\vnu \mid \vxi') p_{\vx \mid \vx_t}(\vxi' \mid \vxi)
  \odif  \vxi'.
\end{equation}
Formally, when the posterior $p_{\vx\mid \vx_t}$ is a delta function centered at
its mean $\bE[\vx \mid \vx_t=\vxi]$, the approximation
\eqref{eq:conditional-posterior-measurementmatching-gaussian-case-dps-approx} is
exact. More generally, when the posterior $p_{\vx \mid \vx_t}$ is highly
concentrated around its mean, the approximation
\eqref{eq:conditional-posterior-measurementmatching-gaussian-case-dps-approx} is
accurate. This holds, for example, for sufficiently small $t$, which we saw
explicitly in the Gaussian setting of
\Cref{example:denoising-conditional-gaussian}.
Although the numerical simulation in
\Cref{fig:conditional_sampling_computational_gaussian} suggests that this
approximation is not without its caveats in certain regimes, it has proved to be
a reliable baseline in practice, after being proposed by Chung et al.\ as
``Diffusion Posterior Sampling'' (DPS) \cite{chung2023diffusion}.
In addition, there are even principled and generalizable approaches to improve it by
incorporating better estimates of the posterior variance (which turn out to be
exact in the Gaussian setting of \Cref{example:denoising-conditional-gaussian}),
which we discuss further in the end-of-chapter summary. 

Thus, with the DPS approximation, we arrive at the following approximation for
the conditional posterior denoisers $\bE[\vx \mid \vy, \vx_t]$, via
\Cref{eq:posterior-sampling-denoiser-decomposition}:
\begin{equation}
  \bE[ \vx \mid \vx_t=\vxi, \vy=\vnu]
  \approx
  \bE[\vx \mid \vx_t=\vxi] 
  + t^2 \nabla_{\vxi} \log p_{\vy \mid \vx}(\vnu \mid \bE[\vx \mid \vx_t = \vxi]).
  \label{eq:posterior-sampling-denoiser-decomposition-dps}
\end{equation}
And, for a neural network or other model $\bar{\vx}_{\theta}(t, \vxi)$ 
trained as in \Cref{sub:compression_denoising} to approximate the denoisers
$\bE[\vx \mid \vx_t = \vxi]$ for each $t \in [0, T]$, we arrive at the learned
conditional posterior denoisers
\begin{equation}
  \bar{\vx}_{\theta}(t, \vxi, \vnu) 
  = \bar{\vx}_{\theta}(t, \vxi)
  + t^2 \nabla_{\vxi} \log p_{\vy \mid \vx}(\vnu \mid \bar{\vx}_{\theta}(t,
  \vxi)).
  \label{eq:posterior-sampling-denoiser-learned-dps}
\end{equation}
Note that the approximation
\eqref{eq:posterior-sampling-denoiser-decomposition-dps} is  valid for arbitrary
forward models $h$ in the
observation model \eqref{eq:measurement-matching-observation}, including
nonlinear $h$, and even to
arbitrary noise models for which a clean expression for the likelihood $p_{\vy
\mid \vx}$ is known. Indeed, in the case of Gaussian noise, we have
\begin{equation}
  p_{\vy \mid \vx}(\vnu \mid \vxi)
  \propto
  \exp\left(
    -\frac{1}{2\sigma^2} \norm*{ h(\vxi) - \vnu }_2^2
  \right).
\end{equation}
Hence, evaluating the righthand side of
\eqref{eq:posterior-sampling-denoiser-learned-dps}
requires only
\begin{enumerate}
  \item A pretrained denoiser $\bar{\vx}_{\theta}(t, \vxi)$ for the data distribution $p$
    (of $\vx)$, learned as in \Cref{sub:compression_denoising} via
    \Cref{alg:learning_denoiser};
  \item Forward and backward pass access to the forward model $h$ for the
    measurements \eqref{eq:measurement-matching-observation};
  \item A forward and backward pass through $\bar{\vx}_{\theta}(t, \vxi)$, which can be evaluated
    efficiently using (say) backpropagation.
\end{enumerate}

\begin{algorithm}
  \caption{Conditional sampling under measurements
  \eqref{eq:measurement-matching-observation}, with an unconditional denoiser and DPS.}
	\label{alg:iterative_denoising_conditional}
	\begin{algorithmic}[1]
		\Require{An ordered list of timesteps \(0 \leq t_{0} < \cdots < t_{L} \leq T\) to use for sampling.}
    \Require{An unconditional denoiser \(\bar{\vx}_{\theta} \colon
    \{t_{\ell}\}_{\ell = 1}^{L} \times \R^{D} \to \R^{D}\) for $p_{\vx}$.}
    \Require{Measurement realization $\vnu$ of $\vy$ (\Cref{eq:measurement-matching-observation}) to condition on.}
    \Require{Forward model $h : \bR^D \to \bR^d$ and measurement noise variance
    $\sigma^2 > 0$.}
		\Require{Scale and noise level functions \(\alpha, \sigma \colon \{t_{\ell}\}_{\ell = 0}^{L} \to \R_{\geq 0}\).}
    \Ensure{A sample \(\hat{\vx}\), approximately from \(p_{\vy \mid \vx}\).}
    \Function{DDIMSamplerConditionalDPS}{$\bar{\vx}_{\theta}, \vnu, h, \sigma^2,
    (t_{\ell})_{\ell = 0}^{L}$}
		\State{Initialize \(\hat{\vx}_{t_{L}} \sim\) approximate distribution of \(\vx_{t_{L}}\)} \Comment{VP \(\implies \dNorm(\vzero, \vI)\), VE \(\implies \dNorm(\vzero, t_{L}^{2}\vI)\).}
		\For{\(\ell = L, L - 1, \dots, 1\)}
		% \State{Compute
		% 	\begin{equation*}
  %       \nabla_{\vxi}\left[ -\frac{1}{2\sigma^2} \norm*{h(
  %       \bar{\vx}_{\theta}(t_{\ell}, \vxi)
  %       ) - \vnu}_2^2 \right]\biggl|_{\vxi = \hat{\vx}_{t_{\ell}}}\biggr.
		% 	\end{equation*}
		% }
		\State{Compute
			\begin{equation*}
				\hat{\vx}_{t_{\ell - 1}} := \frac{\sigma_{t_{\ell
        - 1}}}{\sigma_{t_{\ell}}}\hat{\vx}_{t_{\ell}} + \bp{\alpha_{t_{\ell
        - 1}} - \frac{\sigma_{t_{\ell
        - 1}}}{\sigma_{t_{\ell}}}\alpha_{t_{\ell}}}\left(
        \bar{\vx}_{\theta}(t_{\ell}, \hat{\vx}_{t_{\ell}})
        - \frac{\sigma_{t_{\ell}}^2}{2\alpha_{t_\ell}\sigma^2}
        \nabla_{\vxi}\left[ \norm*{h( \bar{\vx}_{\theta}(t_{\ell}, \vxi)
        ) - \vnu}_2^2 \right]\biggl|_{\vxi = \hat{\vx}_{t_{\ell}}}\biggr.
        \right)
			\end{equation*}
		}
		\EndFor
		\State{\Return{\(\hat{\vx}_{t_{0}}\)}}
		\EndFunction
	\end{algorithmic}
\end{algorithm}

Combining this scheme with the basic implementation of unconditional sampling we
developed in \Cref{sub:compression_denoising}, we obtain a practical algorithm for
conditional sampling of the posterior $p_{\vx \mid \vy}$ given measurements
following \eqref{eq:measurement-matching-observation}.
\Cref{alg:iterative_denoising_conditional} records this scheme
for the case of
Gaussian observation noise with known standard deviation $\sigma$, with minor
modifications to extend to a general noising process, as in
\Cref{eq:gen_additive_gaussian_noise_model} and the surrounding discussion in
\Cref{ch:compression} (our discussion above made the simplifying choices $\alpha_t = 1$, $\sigma_t = t$, and
$t_{\ell} = T\ell / L$, as for \Cref{eq:denoising-iteration-basic} in
\Cref{sub:compression_denoising}).

\subsection{Body Pose Generation Conditioned on Head and Hands} The type of conditional estimation or generation problems arise rather naturally in many practical applications. 
A typical problem of this kind is how to estimate and generate body pose and hand gesture conditioned on a given head pose and egocentric images, as illustrated in Figure \ref{fig:pose-teaser}. This is often the problem we need to solve when one is wearing a head-mounted device such as the Vision-pro from Apple or the Project Aria from Meta. The pose of the whole body and the gesture of the hands need to be inferred so that we can use the information to control virtual objects that the person interacts with.
\begin{figure*}[t]
  \centering
  \includegraphics[width=\linewidth]{figs_chap6/pose-teaser.pdf}
    \caption{
A system that estimates human body height, pose,  and hand parameters (middle), conditioned on ego-centric SLAM poses and images (left). Outputs capture the wearer's actions in the allocentric reference frame of the scene, which we visualize here with 3D reconstructions (right).
  }
  \label{fig:pose-teaser}

\end{figure*}

Notice that in this case, one only has the head pose  provided by the device and a very limited field of view for part of one's hands and upper limps. The pose of the rest of the body needs to ``inferred'' or ``completed'' based on such partial information. The only way one can estimate the body pose at all time is by learning the joint distribution of the head and body pose sequence in advance and then sample this prior distribution conditioned on the real-time partial inputs. Figure \ref{fig:pose-method} outlines a system called EgoAllo \cite{yi2024egoallo} to solve this problem based on a learned conditional diffusion-denoising model. 
\begin{figure*}[t]
  \centering
  \includegraphics[width=\linewidth]{figs_chap6/pose-method.pdf}
\caption{
    \textbf{Overview of technical components of EgoAllo \cite{yi2024egoallo}.}
    % We take SLAM poses and egocentric video as input.
    A diffusion model is pretrained that can generate body pose sequence based on local body parameters (middle).
    An invariant parameterization $g(\cdot)$ of SLAM poses (left) is used to condition the diffusion model. These can be placed into the global coordinate frame via global alignment to input poses.
    When available, egocentric video is used for hand detection (left) via HaMeR~\cite{pavlakos2023reconstructing}, which can be incorporated into samples via guidance by the generated gesture.
  }
  \label{fig:pose-method}
\end{figure*}

Strictly speaking, the solution proposed in EgoAllo \cite{yi2024egoallo} does not enforce measurement matching using the techniques introduced above. Instead it heuristically enforces the condition by utilizing  the cross-attention mechanism in a transformer architecture. There is reason to believe that the cross-attention mechanism is in a way approximately realizing the conditional sampling of the denoising a posteriori. We believe the more principled techniques introduced here, if properly implemented, can lead to better methods that further improve the body pose estimation. 

\section{Conditional Generation with Paired Data}

\subsection{Class Conditioned Image Generation} 
The problem of
classification can be viewed as one to learn a (lossy) encoder for a random vector. That is, given a sample of a random vector $\x$, we would like to predict its class label $\y$. In other words, we need to learn a mapping:
\begin{equation}
  h: \x \mapsto \y.
\end{equation}
In general, we do not know the mapping $h(\cdot)$ and it needs to be learned from a set of paired samples $(\X, \Y) = \{ (\x_1, \y_1), \ldots, (\x_N, \y_N) \}$. 

We know the distribution of $\x$ is low-dimensional. From the previous chapters, we have learned that given sufficient samples, in principle, we can learn a structured representation $\z$ for $\x$ through a learned compressive encoding:
\begin{equation}
    f: \x \mapsto \z. 
\end{equation}
$\z$ can also be viewed as a learned code for $\x$.
It is rather reasonable to assume that the class assignment $\y$ for $\x$ should strongly depend on the (mixed) low-dimensional structures of the distribution, hence those of $\z$ too.

Therefore, we may combine the two codes $\y$ and $\z$ together and try to learn a combined encoder:
\begin{equation}
    f: \x \mapsto (\z, \y). 
\end{equation}
From our study in previous chapters, the mapping $f$ is usually learned as a sequence of compression or denoising operators in the same space. Hence to leverage on such operations, we may introduce another random vector $\vw$ that can be viewed as an initial random guess of the class label $\y$. We try to learn a compression or denoising mapping:
\begin{equation}
    f: (\x, \vw) \mapsto (\z, \y)
\end{equation}
within a common feature space. 

In fact, the common practice of introducing an
auxiliary ``class token'' in the training of a Transformer for
classification tasks, such as in ViT, can be viewed as learning such a
representation by compressing (the coding rate of) given (noisy) samples of
$(\x, \vw)$. If the distribution of the data $\x$ is already a
mixture of (low-dimensional) Gaussians, the work \cite{wright2008classification} has shown that classification can be
done effectively by directly minimizing {\em the (lossy) coding
length} associated with the given samples.


\subsection{Caption Conditioned Image Generation}
An vignette for image generation conditioned on a given caption. (say, DALLE with CLIP for caption conditioned generation.)



\section{Conditional Generation with Measurement Self-Consistency}
In this last section, we consider the more extreme case in which we only have a set of observed samples $\Y = \{\y_1,\ldots, \y_N\}$ of the data $\x$. To make the problem well-defined, we do assume that the observation model $h(\cdot)$ between $\y$ and $\x$ is relatively simple and known. 

\subsection{New View Synthesis from a Sequence of Images}
Learn to render new views of a static 3D scene from a sequence of uncalibrated (monocular) images. Or learn to render new views of a dynamical 3.5D scene from a sequence of uncalibrated stereo image pairs. 


\section{Summary and Notes}
Allowing the prior distribution to have a finite low-dimensional support helps reconcile the gap between probabilistic inference and deterministic inference, based on whether the given condition is sufficient to make the conditional estimation or generation deterministic. 

\section{Exercises}

\end{document}